% ============================================================
\chapter{Introduction au MLOps --- Du Notebook \`a la Production}
\label{ch:introduction}
\index{MLOps!introduction}
% ============================================================

\section{Qu'est-ce que le MLOps~?}
\label{sec:mlops-definition}

Le \textbf{MLOps}\index{MLOps} d\'esigne l'ensemble des pratiques, outils et
principes culturels qui visent \`a d\'eployer et maintenir des mod\`eles de
machine learning en production de mani\`ere fiable et efficace. Le terme est
un mot-valise combinant \emph{Machine Learning} et \emph{Operations}, par
analogie avec le \emph{DevOps}\index{DevOps} du g\'enie logiciel.

\begin{definition}[MLOps]
Le MLOps est la discipline qui applique les principes du DevOps (int\'egration
continue, d\'eploiement continu, monitoring, automatisation) au cycle de vie
des mod\`eles de machine learning, en y ajoutant les sp\'ecificit\'es li\'ees
aux donn\'ees, aux exp\'eriences et \`a la d\'erive des mod\`eles.
\end{definition}

\subsection{DevOps vs MLOps}

Le DevOps traditionnel g\`ere du code. Le MLOps g\`ere du code \textbf{et}
des donn\'ees \textbf{et} des mod\`eles --- trois artefacts dont l'interaction
cr\'ee une complexit\'e suppl\'ementaire.

\begin{center}
\begin{tikzpicture}[
  box/.style={rectangle, draw, rounded corners, minimum width=3cm,
    minimum height=1.5cm, align=center, font=\small},
  arr/.style={-{Stealth[length=2.5mm]}, thick}
]
  % DevOps
  \node[box, fill=blue!10] (devops) {\textbf{DevOps}\\Code};
  % MLOps
  \node[box, fill=green!10, right=3cm of devops] (mlops) {\textbf{MLOps}\\Code + Donn\'ees\\+ Mod\`eles};

  % Shared
  \node[above=1.5cm of devops, font=\small\bfseries] (shared) {Pratiques partag\'ees};
  \node[below=0.3cm of shared, font=\scriptsize, align=center] {CI/CD, conteneurisation,\\monitoring, IaC};

  \draw[arr, steelblue] (shared.south) -- (devops.north);
  \draw[arr, steelblue] (shared.south) -- (mlops.north);

  % MLOps specific
  \node[below=0.8cm of mlops, font=\scriptsize, align=center, text=deepred] {+ Versionnement donn\'ees\\+ Suivi exp\'eriences\\+ D\'etection de drift\\+ Feature engineering};
\end{tikzpicture}
\end{center}

\begin{remarque}
La diff\'erence fondamentale est que le comportement d'un syst\`eme ML
d\'epend non seulement du code mais aussi des donn\'ees. Un m\^eme code
entra\^in\'e sur des donn\'ees diff\'erentes produit un mod\`ele diff\'erent.
Cela n\'ecessite de versionner les donn\'ees au m\^eme titre que le code.
\end{remarque}

\section{Les niveaux de maturit\'e MLOps}
\label{sec:maturity-levels}
\index{MLOps!niveaux de maturit\'e}

Google a propos\'e trois niveaux de maturit\'e pour les syst\`emes ML en
production~:

\begin{center}
\begin{tabular}{@{}llp{7cm}@{}}
\toprule
\textbf{Niveau} & \textbf{Nom} & \textbf{Description} \\
\midrule
0 & Manuel & Processus enti\`erement manuel. Notebooks, pas de pipeline.
  S\'eparation compl\`ete entre data scientists et ops. \\
1 & Pipeline ML & Pipeline automatis\'e d'entra\^inement. D\'eploiement
  continu du mod\`ele. Monitoring basique. \\
2 & Pipeline CI/CD & Pipeline CI/CD pour le code ML. Entra\^inement
  automatis\'e d\'eclench\'e par les donn\'ees. R\'e-entra\^inement
  automatique en cas de drift. \\
\bottomrule
\end{tabular}
\end{center}

\begin{center}
\begin{tikzpicture}[
  level/.style={rectangle, draw, rounded corners, minimum width=10cm,
    minimum height=1.2cm, align=center, font=\small},
  arr/.style={-{Stealth[length=2.5mm]}, thick, steelblue}
]
  \node[level, fill=red!10] (l0) {\textbf{Niveau 0~: Manuel}\\
    Notebooks $\to$ export mod\`ele $\to$ d\'eploiement manuel};
  \node[level, fill=yellow!10, above=0.5cm of l0] (l1) {\textbf{Niveau 1~: Pipeline ML}\\
    Pipeline automatis\'e $\to$ d\'eploiement continu $\to$ monitoring};
  \node[level, fill=green!10, above=0.5cm of l1] (l2) {\textbf{Niveau 2~: CI/CD + ML}\\
    CI/CD code + donn\'ees $\to$ r\'e-entra\^inement auto $\to$ A/B testing};

  \draw[arr] (l0) -- (l1);
  \draw[arr] (l1) -- (l2);
\end{tikzpicture}
\end{center}

\section{Le cycle de vie ML}
\label{sec:ml-lifecycle}
\index{cycle de vie ML}

Un projet ML en production traverse plusieurs phases, qui forment un cycle
(et non une s\'equence lin\'eaire)~:

\begin{enumerate}
  \item \textbf{D\'efinition du probl\`eme}~: quel est l'objectif m\'etier~?
    Quel est le crit\`ere de succ\`es~?
  \item \textbf{Collecte et pr\'eparation des donn\'ees}~: acquisition,
    nettoyage, labellisation, feature engineering.
  \item \textbf{Exp\'erimentation}~: choix du mod\`ele, optimisation des
    hyperparam\`etres, \'evaluation.
  \item \textbf{D\'eveloppement du pipeline}~: transformation du notebook
    en code modulaire, test\'e et versionn\'e.
  \item \textbf{D\'eploiement}~: mise en production du mod\`ele (API,
    batch, edge).
  \item \textbf{Monitoring}~: surveillance de la performance, d\'etection
    de drift, alertes.
  \item \textbf{R\'e-entra\^inement}~: mise \`a jour du mod\`ele quand les
    donn\'ees ou les performances \'evoluent.
\end{enumerate}

\begin{center}
\begin{tikzpicture}[
  node distance=1.8cm,
  phase/.style={circle, draw, minimum size=1.6cm, align=center,
    font=\scriptsize\bfseries, inner sep=1pt},
  arr/.style={-{Stealth[length=2mm]}, thick, steelblue}
]
  \node[phase, fill=blue!15] (prob) {Probl\`eme};
  \node[phase, fill=blue!10, right=1.5cm of prob] (data) {Donn\'ees};
  \node[phase, fill=green!10, below right=1cm and 1cm of data] (exp) {Exp\'erim.};
  \node[phase, fill=green!15, below=1.5cm of exp] (pipe) {Pipeline};
  \node[phase, fill=orange!10, below left=1cm and 1cm of pipe] (deploy) {D\'eploie.};
  \node[phase, fill=red!10, left=1.5cm of deploy] (monitor) {Monitoring};
  \node[phase, fill=red!15, above left=1cm and 1cm of monitor] (retrain) {R\'e-entra\^in.};

  \draw[arr] (prob) -- (data);
  \draw[arr] (data) -- (exp);
  \draw[arr] (exp) -- (pipe);
  \draw[arr] (pipe) -- (deploy);
  \draw[arr] (deploy) -- (monitor);
  \draw[arr] (monitor) -- (retrain);
  \draw[arr] (retrain) -- (data);
\end{tikzpicture}
\end{center}

\section{La dette technique en ML}
\label{sec:technical-debt}
\index{dette technique}

Sculley et al.~\cite{sculley2015} ont identifi\'e que le code ML ne
repr\'esente qu'une petite fraction d'un syst\`eme ML en production. Le reste
est de l'infrastructure~:

\begin{center}
\begin{tikzpicture}[
  box/.style={rectangle, draw, minimum height=0.8cm, align=center,
    font=\scriptsize}
]
  % Central ML code
  \node[box, fill=black!80, text=white, minimum width=2cm] (ml) {Code ML};

  % Surrounding infrastructure
  \node[box, fill=yellow!20, minimum width=3.5cm, above=0.1cm of ml] (config) {Configuration};
  \node[box, fill=yellow!20, minimum width=3.5cm, below=0.1cm of ml] (data) {Collecte de donn\'ees};
  \node[box, fill=yellow!20, minimum width=1.5cm, left=0.1cm of ml] (feat) {Feature\\engineering};
  \node[box, fill=yellow!20, minimum width=1.5cm, right=0.1cm of ml] (serve) {Infrastructure\\de serving};
  \node[box, fill=orange!20, minimum width=6.5cm, above=0.1cm of config] (monitor) {Monitoring \quad Outils de processus \quad Gestion des ressources};
  \node[box, fill=orange!20, minimum width=6.5cm, below=0.1cm of data] (valid) {V\'erification donn\'ees \quad Analyse \quad Tests};
\end{tikzpicture}
\end{center}

\begin{attention}[title=\textbf{Le pi\`ege du notebook}]
Le notebook Jupyter est un excellent outil d'exploration, mais un pi\`etre
outil de production. Les probl\`emes classiques~:
\begin{itemize}
  \item Ex\'ecution non lin\'eaire (cellules ex\'ecut\'ees dans le d\'esordre)
  \item \'Etat cach\'e (variables globales qui persistent)
  \item Pas de tests unitaires
  \item Difficult\'e de versionnement (fichiers JSON volumineux)
  \item Pas de modularit\'e (fonctions et classes non r\'eutilisables)
\end{itemize}
\end{attention}

\section{Architecture d'un syst\`eme MLOps}
\label{sec:mlops-architecture}

Un syst\`eme MLOps mature comprend les composants suivants~:

\begin{toolbox}[title=\textbf{Composants d'un syst\`eme MLOps}]
\begin{description}
  \item[Gestion de code] Git, GitHub/GitLab
  \item[Gestion de donn\'ees] DVC, Delta Lake, LakeFS
  \item[Gestion d'environnements] Conda, pip, venv, Poetry
  \item[Suivi d'exp\'eriences] MLflow, Weights \& Biases, Neptune
  \item[Pipeline d'entra\^inement] Airflow, Prefect, Kubeflow Pipelines
  \item[Feature store] Feast, Tecton, Hopsworks
  \item[Conteneurisation] Docker, Kubernetes
  \item[Serving] FastAPI, TorchServe, TensorFlow Serving, Triton
  \item[CI/CD] GitHub Actions, GitLab CI, Jenkins
  \item[Monitoring] Prometheus, Grafana, Evidently, WhyLabs
\end{description}
\end{toolbox}

\section{Tutoriel~: du notebook au script modulaire}
\label{sec:notebook-to-script}

\subsection{Le notebook de d\'epart}

Consid\'erons un notebook typique de classification~:

\begin{codeblock}[title=\textbf{Notebook typique (anti-pattern)}]
\begin{minted}{python}
# Cellule 1
import pandas as pd
from sklearn.ensemble import RandomForestClassifier
from sklearn.model_selection import train_test_split
from sklearn.metrics import accuracy_score

# Cellule 2
df = pd.read_csv("data.csv")
X = df.drop("target", axis=1)
y = df["target"]

# Cellule 3
X_train, X_test, y_train, y_test = train_test_split(
    X, y, test_size=0.2, random_state=42
)

# Cellule 4
model = RandomForestClassifier(n_estimators=100, random_state=42)
model.fit(X_train, y_train)
print(f"Accuracy: {accuracy_score(y_test, model.predict(X_test)):.4f}")
\end{minted}
\end{codeblock}

\subsection{Le script modulaire}

\begin{bonnepratique}[title=\textbf{Structure de projet recommand\'ee}]
\begin{minted}{text}
my_ml_project/
  src/
    __init__.py
    data.py          # Chargement et preprocessing
    model.py         # Definition du modele
    train.py         # Boucle d'entrainement
    evaluate.py      # Metriques d'evaluation
  tests/
    test_data.py
    test_model.py
  configs/
    config.yaml      # Hyperparametres
  requirements.txt
  Makefile
  README.md
\end{minted}
\end{bonnepratique}

\begin{codeblock}[title=\textbf{src/data.py --- Module de donn\'ees}]
\begin{minted}{python}
"""Data loading and preprocessing module."""
import pandas as pd
from sklearn.model_selection import train_test_split


def load_data(path: str) -> pd.DataFrame:
    """Load dataset from CSV file."""
    df = pd.read_csv(path)
    return df


def split_data(
    df: pd.DataFrame,
    target_col: str = "target",
    test_size: float = 0.2,
    random_state: int = 42,
) -> tuple:
    """Split data into train and test sets."""
    X = df.drop(target_col, axis=1)
    y = df[target_col]
    return train_test_split(
        X, y, test_size=test_size, random_state=random_state
    )
\end{minted}
\end{codeblock}

\begin{codeblock}[title=\textbf{src/train.py --- Script d'entra\^inement}]
\begin{minted}{python}
"""Training script with configuration."""
import argparse
import yaml
import joblib
from data import load_data, split_data
from sklearn.ensemble import RandomForestClassifier
from sklearn.metrics import accuracy_score


def train(config: dict) -> None:
    """Train model with given configuration."""
    # Data
    df = load_data(config["data"]["path"])
    X_train, X_test, y_train, y_test = split_data(
        df,
        target_col=config["data"]["target_col"],
        test_size=config["data"]["test_size"],
    )

    # Model
    model = RandomForestClassifier(**config["model"])
    model.fit(X_train, y_train)

    # Evaluate
    y_pred = model.predict(X_test)
    acc = accuracy_score(y_test, y_pred)
    print(f"Accuracy: {acc:.4f}")

    # Save
    joblib.dump(model, config["output"]["model_path"])
    print(f"Model saved to {config['output']['model_path']}")


if __name__ == "__main__":
    parser = argparse.ArgumentParser()
    parser.add_argument("--config", default="configs/config.yaml")
    args = parser.parse_args()

    with open(args.config) as f:
        config = yaml.safe_load(f)

    train(config)
\end{minted}
\end{codeblock}

\begin{codeblock}[title=\textbf{configs/config.yaml}]
\begin{minted}{yaml}
data:
  path: "data/dataset.csv"
  target_col: "target"
  test_size: 0.2

model:
  n_estimators: 100
  max_depth: 10
  random_state: 42

output:
  model_path: "models/rf_model.joblib"
\end{minted}
\end{codeblock}

\begin{codeblock}[title=\textbf{Makefile}]
\begin{minted}{makefile}
.PHONY: install train test clean

install:
	pip install -r requirements.txt

train:
	python src/train.py --config configs/config.yaml

test:
	pytest tests/ -v

clean:
	rm -rf models/*.joblib __pycache__ .pytest_cache
\end{minted}
\end{codeblock}

\section{Bonnes pratiques et anti-patterns}

\begin{bonnepratique}[title=\textbf{Principes fondamentaux}]
\begin{enumerate}
  \item \textbf{S\'eparation des pr\'eoccupations}~: donn\'ees, mod\`ele,
    entra\^inement, \'evaluation dans des modules distincts.
  \item \textbf{Configuration externalis\'ee}~: hyperparamm\`etres dans des
    fichiers YAML, pas en dur dans le code.
  \item \textbf{Reproductibilit\'e par d\'efaut}~: fixer les graines
    al\'eatoires, versionner les donn\'ees et le code.
  \item \textbf{Tests automatiques}~: tester le chargement des donn\'ees,
    la forme des tenseurs, la convergence sur un petit \'echantillon.
  \item \textbf{Automatisation}~: Makefile ou scripts pour toutes les
    op\'erations r\'ep\'etitives.
\end{enumerate}
\end{bonnepratique}

\begin{antipattern}[title=\textbf{\`A \'eviter absolument}]
\begin{itemize}
  \item Hardcoder des chemins absolus~: \cli{/home/alice/data/train.csv}
  \item Committer des donn\'ees volumineuses dans Git
  \item Committer des secrets (cl\'es API, mots de passe)
  \item Travailler exclusivement en notebooks sans code modulaire
  \item Ignorer les tests sous pr\'etexte que <<~c'est du ML~>>
  \item D\'eployer un mod\`ele sans monitoring
\end{itemize}
\end{antipattern}

\section{Mini-projet~: structurer un projet ML}
\label{sec:mini-project-ch1}

\begin{miniprojet}[title=\textbf{Mini-projet 1~: Du notebook au projet structur\'e}]
\begin{enumerate}
  \item Cr\'eez la structure de r\'epertoires recommand\'ee.
  \item Transformez le notebook de classification ci-dessus en modules
    Python s\'epar\'es.
  \item \'Ecrivez un fichier \file{config.yaml} pour les hyperparamm\`etres.
  \item Cr\'eez un \file{Makefile} avec les cibles \cli{install},
    \cli{train}, \cli{test}, \cli{clean}.
  \item \'Ecrivez au moins deux tests unitaires (un pour le chargement des
    donn\'ees, un pour la forme de la sortie du mod\`ele).
  \item V\'erifiez que \cli{make train} fonctionne de bout en bout.
\end{enumerate}
\end{miniprojet}

\section{Exercices}
\label{sec:exercises-ch1}

\begin{exercice}[$\bigstar$ --- Structure de projet]
Cr\'eez un squelette de projet ML pour un probl\`eme de r\'egression avec
la structure recommand\'ee. Incluez un \file{README.md} d\'ecrivant le
projet et un \file{.gitignore} adapt\'e au ML (ignorer \file{data/},
\file{models/}, \file{*.pyc}, \file{.ipynb\_checkpoints/}).
\end{exercice}

\begin{exercice}[$\bigstar\bigstar$ --- Refactoring de notebook]
Prenez un notebook Jupyter existant (le v\^otre ou un notebook Kaggle) et
transformez-le en projet structur\'e~:
\begin{enumerate}
  \item Identifiez les fonctions r\'eutilisables
  \item Cr\'eez les modules Python correspondants
  \item Externalisez la configuration
  \item \'Ecrivez des tests pour les fonctions critiques
\end{enumerate}
\end{exercice}

\begin{exercice}[$\bigstar\bigstar\bigstar$ --- Projet complet]
Construisez un projet ML complet pour la classification de texte
(par exemple, analyse de sentiment sur le dataset IMDb)~:
\begin{enumerate}
  \item Structure modulaire compl\`ete
  \item Configuration YAML
  \item Tests unitaires avec \pkg{pytest}
  \item Script d'entra\^inement avec arguments en ligne de commande
  \item Rapport automatique des m\'etriques (pr\'ecision, rappel, F1)
  \item \file{Makefile} complet
\end{enumerate}
\end{exercice}

\section{Aide-m\'emoire}

\begin{cheatsheet}[title=\textbf{R\'esum\'e du chapitre 1}]
\begin{tabular}{@{}p{4cm}p{8.5cm}@{}}
\toprule
\textbf{Concept} & \textbf{Point cl\'e} \\
\midrule
MLOps & DevOps + donn\'ees + mod\`eles \\
Niveaux de maturit\'e & 0 (manuel) $\to$ 1 (pipeline) $\to$ 2 (CI/CD) \\
Dette technique & Le code ML est une petite fraction du syst\`eme \\
Anti-pattern notebook & \'Etat cach\'e, pas de tests, pas de modularit\'e \\
Structure projet & \file{src/}, \file{tests/}, \file{configs/}, Makefile \\
Config externalis\'ee & YAML pour les hyperparamm\`etres \\
\bottomrule
\end{tabular}
\end{cheatsheet}
